% =========================================================================
% config.tex -- master configuration for the BOOK layout
% -------------------------------------------------------------------------
% Same package family and CUSTOM-COMMAND NAMES as config_manuscript.tex, so
% a manuscript chapter needs only its call PARAMETERS updated (done by
% transform_manuscript.py), not its structure.
%
% All tunable numbers/colours/toggles live in parameters.tex (\input below).
% This file is the workhorse: packages, macros, environments, page logic.
% Section numbers match parameters.tex.
% =========================================================================

% -----------------------------------------------------------------------
% Packages
% -----------------------------------------------------------------------
\usepackage[utf8]{inputenc}
\usepackage[T1]{fontenc}
\usepackage{amsmath, amssymb}
\usepackage{graphicx}
\usepackage{booktabs}
\usepackage{threeparttable}
\usepackage{array}
\usepackage[english]{babel}
\usepackage{hyperref}   % must load before cleveref
\usepackage{cleveref}
\usepackage{xstring}
\usepackage{geometry}
\usepackage{mdframed}
\usepackage{tcolorbox}
  \tcbuselibrary{skins, breakable}
\usepackage{tikz}       % needed for \bookcropimage (bio portrait centre-crop)
\usepackage{xcolor}
\newif\ifbooksoulloaded
\IfFileExists{soul.sty}{\usepackage{soul}\booksoulloadedtrue}{\booksoulloadedfalse}
\ifbooksoulloaded\else\providecommand{\hl}[1]{\textcolor{red}{#1}}\fi
\usepackage{setspace}
\usepackage{microtype}
\usepackage[shortlabels]{enumitem}
\usepackage{ragged2e}
\usepackage{helvet}
\usepackage{times}
\usepackage{titlesec}
\usepackage{eurosym}
\usepackage{changepage}
\usepackage{wasysym}
\usepackage{fancyhdr}
\usepackage{lastpage}
\usepackage{eso-pic}    % full-bleed cover image
\usepackage{multicol}   % balanced spoiler columns on the chapter cover
\IfFileExists{lipsum.sty}{\usepackage{lipsum}}{}
\usepackage{stfloats}   % required for [b] placement on figure*/table* floats
\usepackage{lineno}     % editorial line numbers (off by default, see parameters.tex)
\usepackage[font={footnotesize,sf}, labelfont={bf,color=headingcolor}, labelsep=quad]{caption}
\usepackage[useregional]{datetime2}   % needed for \DTMNow on the colophon page
\usepackage{etoolbox}

% -----------------------------------------------------------------------
% Toggles & base setup (declared here; state set in parameters.tex)
% -----------------------------------------------------------------------
\newif\ifchapterstartseven   % which parity chapters open on
\newif\ifbooklinenumbers     % editorial line numbers on/off
\newif\ifbookboxsizeS        % internal: tracks S vs L for \boxmain, not user-facing

\renewcommand{\linenumberfont}{\normalfont\tiny\sffamily\color{gray}}
\renewcommand{\sfdefault}{phv}  % phv = Helvetica

\newcommand{\degree}{\ensuremath{{}^{\circ}}}
\newcommand{\authornote}[1]{\hl{[Author note: #1]}}

\newlength{\bookvunit}
\AtBeginDocument{\setlength{\bookvunit}{\baselineskip}}

% -----------------------------------------------------------------------
% Load all tunable values
% -----------------------------------------------------------------------
% =========================================================================
% parameters.tex -- USER-TUNABLE VALUES ONLY.
%
% This file is the single place to edit numbers, colours, and toggles for
% the book layout. It is \input by config.tex, which reads these macros
% to build the actual page geometry, styles, and environments.
%
% RULE OF THUMB: if you're tempted to open config.tex to "just change a
% number", that number should probably live here instead. config.tex
% should only need edits when you're changing LOGIC (how something is
% built), not VALUES (how big/far/what colour something is).
%
% Section numbers below match the corresponding sections in config.tex,
% so "3. LINES, PARAGRAPHS..." here lines up with the same heading there.
% Some numbers (4, 8, 11...) are skipped on purpose -- those sections in
% config.tex are logic-only and have nothing to tune here.
% =========================================================================

% -----------------------------------------------------------------------
% 1. PAGE GEOMETRY
% -----------------------------------------------------------------------
% The physical page and the two-column text block that sits on it.
% Everything here is in absolute units (mm) so the numbers are easy to
% check against a printed spec sheet.

\newcommand{\bookpagewidth}{216mm}      % trim width
\newcommand{\bookpageheight}{276mm}     % trim height

\newcommand{\bookcolwidth}{83.5mm}      % width of ONE text column
\newcommand{\bookcolsep}{5mm}           % gap between the two columns
\newcommand{\booktextwidth}{172mm}      % MUST equal 2*\bookcolwidth + \bookcolsep
                                         % (not computed automatically -- if you
                                         % change colwidth/colsep, update this too)

\newcommand{\booklatmargin}{22mm}       % left/right margin, symmetric.
                                         % Should equal (\bookpagewidth - \booktextwidth)/2

% --- Vertical space budget, top of page -------------------------------
% The header lives INSIDE the top margin (geometry package "includehead").
% Reading top to bottom, the top margin is built from three stacked pieces:
\newcommand{\booktopblank}{5mm}         % blank strip above the header text
\newcommand{\bookheadheight}{7mm}       % box that holds the header text itself
\newcommand{\bookheadsep}{3mm}          % gap between header rule and first body line
% Sum of the three above = 15mm = actual top margin as far as `geometry` is
% concerned. That is NOT the same variable as the one below, which is used
% independently by the cover-image cropping math (Section 13). If you
% change the three values above, you likely also need to update
% \booktopmargintotal to keep the cover math visually consistent.
% >>> KNOWN OPEN ISSUE: these currently disagree (15mm vs 25mm). <
\newcommand{\booktopmargintotal}{25mm}  % top margin as used by cover-image math only

\newcommand{\bookbottommargin}{20mm}    % blank strip below body text (no footer used)
\newcommand{\bookfootskip}{1pt}         % geometry footskip -- kept minimal since
                                         % there is no footer in this design


% -----------------------------------------------------------------------
% 2. COLOURS
% -----------------------------------------------------------------------
% Core palette. Anything that isn't part-dependent lives at the top;
% part-dependent box colours (Part 1/2/3) are grouped below with the
% chapter-range switches that decide which colour applies where.

\definecolor{maintextcolor}{RGB}{33,30,31}      % body text (near-black, not pure black)
\definecolor{headingcolor}{RGB}{42,110,184}     % headings, header rule, emphasis words
\definecolor{boxcolor}{RGB}{244,247,252}        % fallback/default box fill
\definecolor{bookcovernumcolor}{RGB}{205,205,205} % chapter-number digit on cover bar
\definecolor{bookcovercapcolor}{RGB}{111,111,111} % cover caption text colour

% --- Part-dependent box colour ------------------------------------------
% The book is split into three Parts. Boxes (tips, examples, etc.) pick up
% a tint depending on which Part the current chapter belongs to.
\definecolor{boxcolorpart1}{RGB}{244,247,252}   % Part 1 -- light blue
\definecolor{boxcolorpart2}{RGB}{238,252,250}   % Part 2 -- light green
\definecolor{boxcolorpart3}{RGB}{252,246,238}   % Part 3 -- light amber

% Part boundaries, given as "first chapter number NOT in this Part".
% e.g. with the values below: Part 1 = chapters 0-5, Part 2 = 6-9,
% Part 3 = 10 onward. Edit these two numbers to move the Part breaks.
\newcommand{\bookpartoneend}{6}   % first chapter OUTSIDE Part 1
\newcommand{\bookparttwoend}{10}  % first chapter OUTSIDE Part 2

\colorlet{headerrulecolor}{headingcolor}  % the horizontal rule under the header
\newcommand{\headerrulethickness}{1.5pt}  % rule thickness, spans full text width


% -----------------------------------------------------------------------
% 3. LINES, PARAGRAPHS, HEADINGS & ITEMS SPACING
% -----------------------------------------------------------------------
% Vertical rhythm of the running text. Values are given explicitly in mm
% (not auto-derived from \baselineskip) -- an earlier "grid-snap to
% \bookvunit" approach was tried and dropped as unnecessary complexity;
% explicit before/after values are simpler to reason about and tune.

% --- Global text shape ---------------------------------------------------
\newcommand{\booklinespread}{1.0}   % line-spacing multiplier (1.0 = single)
\newcommand{\bookparindent}{5mm}    % first-line paragraph indent
\newcommand{\bookparskip}{2mm}      % vertical gap left BETWEEN paragraphs

% --- Heading spacing -------------------------------------------------
% Each heading level has an independent "space before" / "space after".
% "Before" = gap above the heading (end of previous content -> heading).
% "After"  = gap below the heading (heading -> start of its content).

% \section
\newcommand{\bookspacebeforesection}{3.0mm}
\newcommand{\bookspaceaftersection}{1.0mm}

% \subsection
\newcommand{\bookspacebeforesubsection}{2.0mm}
\newcommand{\bookspaceaftersubsection}{1.0mm}

% \subsubsection
\newcommand{\bookspacebeforesubsubsection}{1.2mm}
\newcommand{\bookspaceaftersubsubsection}{0.0mm}

% \paragraph (run-in heading, stays on the same line as its text)
\newcommand{\bookspacebeforeparagraph}{1.0mm}
\newcommand{\bookspaceafterparagraph}{5mm}   % HORIZONTAL gap: run-in label -> text
                                              % (different from the others, which
                                              % are all vertical gaps)

% Worked-example blocks
\newcommand{\bookspacebeforeexample}{5.0mm}

% --- Section numbering ----------------------------------------------
\newcommand{\booksecnumdepth}{3}          % how deep numbering goes (3 = down to \subsubsection)
\newcommand{\booksecnumwidth}{10mm}       % label box width for "12"       (\thesection)
\newcommand{\booksubsecnumwidth}{10mm}    % label box width for "12.10"    (\thesubsection)
\newcommand{\booksubsubsecnumwidth}{10mm} % label box width for "(a)"      (\thesubsubsection)
% If a number ever looks clipped or oddly spaced, check these widths first --
% they're fixed boxes, not auto-sized to the label text.

% --- Float spacing --------------------------------------------------
\newcommand{\bookfloatsep}{5mm}                 % gap BETWEEN two stacked floats
\newcommand{\booktextfloatsep}{0.75\bookvunit}  % gap between a float and body text
\newcommand{\bookintextsep}{1.0\bookvunit}      % gap around an in-text (non-floating) figure

% --- Display-equation spacing -----------------------------------------
% Applies to equation, align, subequations, etc.
\newlength{\bookeqspaceabove}  \setlength{\bookeqspaceabove}{0pt}
\newlength{\bookeqspacebelow}  \setlength{\bookeqspacebelow}{2pt}

% --- Table of Contents -------------------------------------------------
\newcommand{\bookcontentstitlesize}{\LARGE}          % size of the "Contents" heading
\newcommand{\bookspacebeforechapter}{0pt}            % gap above "Contents" title
\newcommand{\bookspaceafterchapter}{1.0\bookvunit}   % gap: "Contents" title -> first entry


% -----------------------------------------------------------------------
% 5. CHAPTER OPENING / EDITORIAL TOGGLES
% -----------------------------------------------------------------------
% Booleans (\...true / \...false) rather than numbers -- flip these to
% change document-wide behaviour without touching config.tex.

\chapterstartseventrue    % chapters open on an EVEN page.
                          % Use \chapterstartsevenfalse to force ODD instead.
\booklinenumbersfalse     % editorial line numbers (for review passes) -- OFF by
                          % default. Flip to \booklinenumberstrue for a review draft.

\newcommand{\bookfirstdecoratedchapter}{1}
% Chapters with number < this value are treated as "undecorated" -- they get
% the plain chapter anatomy (no cover spoiler, no skills/take-home boxes).
% This is how Chapter 0 (Preface) gets its stripped-down layout: it's simply
% below \bookfirstdecoratedchapter. Raise this number if you add more
% front-matter-style chapters that shouldn't get the full decoration.


% -----------------------------------------------------------------------
% 6. CHAPTER COVER PAGE
% -----------------------------------------------------------------------
% The full-bleed opening page of each (decorated) chapter: photo at top,
% a number/caption bar, then a two-column "in this chapter" spoiler text.
% Reading order below follows the page top-to-bottom.

\newcommand{\bookcoverimageheight}{160mm}  % height of the full-bleed photo, from page top
\newcommand{\bookcoverbarheight}{32mm}     % MINIMUM height of the number/caption bar
                                            % (grows if the caption text needs more room)
\newcommand{\bookcovernumwidth}{32mm}      % width of the chapter-number box
                                            % (kept close to bar height for a square look)
\newcommand{\bookcovernumfont}{\fontsize{70}{75}\selectfont}  % size of the number digit(s)
\newcommand{\bookcovercapwidth}{137mm}     % width of the caption text area (bar minus number box)
\newcommand{\bookcoverbarpad}{4mm}         % padding: caption text -> bar's minimum height
\newcommand{\bookcovergaptobar}{4mm}       % gap: bottom of image -> top of bar
\newcommand{\bookcoverlabelgap}{3mm}       % gap: "In this chapter..." label -> spoiler columns
\newcommand{\bookcoverspoilerheight}{70mm} % FIXED total height of the spoiler block
                                            % (this is a hard limit, not a minimum --
                                            % trim the spoiler wording if text overflows it)
\newcommand{\bookcovercapspread}{1.3}      % line-spacing multiplier, caption text only


% -----------------------------------------------------------------------
% 7. CHAPTER FIRST PAGE
% -----------------------------------------------------------------------
% The page right after the cover: chapter title + number box on top,
% followed by the two-column body (with the skills box in column 2, top).

\newcommand{\bookfirstpagetitlesize}{\fontsize{24}{26}\selectfont}
\newcommand{\bookfirstpagetitlegap}{10mm}      % gap: title row -> start of column 1 body
\newcommand{\bookfirstpagetitlenumgap}{4mm}    % gap: title text block -> number box
\newcommand{\bookfirstpagetitlenudge}{0.15mm}  % small vertical correction so the title
                                                % row optically aligns with \bookheadsep --
                                                % only touch this if the title looks
                                                % slightly off vs. the header rule
\newcommand{\bookfirstpagenumboxwidth}{30mm}
\newcommand{\bookfirstpagenumboxheight}{35mm}


% -----------------------------------------------------------------------
% 9. FIGURES
% -----------------------------------------------------------------------
% Three figure sizes, selected via the S / M / L flag passed to the figure
% command. S and L placeholders scale to a FIXED height (real images only
% scale by width, keeping their own aspect ratio); M always uses the fixed
% \bookMfigwidth with a side caption, handled as its own minipage layout.

% --- Widths ---------------------------------------------------------
\newcommand{\bookfigwidthS}{\linewidth}   % S = one column wide
\newcommand{\bookMfigwidth}{126mm}        % M = fixed width, side caption
\newcommand{\bookfigwidthL}{\linewidth}   % L = two columns wide (= \booktextwidth)

% --- Placeholder heights ----------------------------------------------
% Only used when drawing a placeholder box (no real image yet). Real
% images ignore these and scale purely by width.
\newcommand{\bookfigheightS}{70mm}
\newcommand{\bookfigheightM}{100mm}
\newcommand{\bookfigheightL}{120mm}

% --- Caption spacing (S and L only; M is handled by its own minipage) ---
\newcommand{\bookfigcapabove}{0.8\bookvunit}   % gap: image -> caption
\newcommand{\bookfigcapbelow}{0.2\bookvunit}   % gap: caption -> content that follows

% --- Placeholder labels ------------------------------------------------
% Text shown inside a placeholder box so you can tell sizes apart at a
% glance while drafting, before real artwork is dropped in.
\newcommand{\bookfiglabelS}{S \textbar{} 83.5 $\times$ 70 mm}
\newcommand{\bookfiglabelM}{M \textbar{} 126 $\times$ 100 mm}
\newcommand{\bookfiglabelL}{L \textbar{} 172 $\times$ 120 mm}

% --- Special-purpose figures --------------------------------------------
\newcommand{\bookboxfigheight}{45mm}    % placeholder height for figures INSIDE boxes
\newcommand{\bookbiofigheight}{25mm}    % author bio portrait -- fixed crop height
\newcommand{\bookbiofigwidth}{18.75mm}  % bio portrait width (3:4 ratio => 0.75 * height)
\newcommand{\bookbiotitlegap}{3mm}      % gap: bio name/dates -> body text
\newcommand{\bookbiogap}{2mm}           % gap: bio image pane -> bio text pane


% -----------------------------------------------------------------------
% 10. BOXES
% -----------------------------------------------------------------------
% Inner padding for tcolorbox-based boxes (tips, examples, take-home,
% skills, etc.). All measured from the box border inward.

\newcommand{\bookboxpad}{1mm}             % tcolorbox "boxsep" -- outer padding, all sides
\newcommand{\bookboxpadlat}{0.5mm}        % extra left/right inner padding (body text)
\newcommand{\bookboxpadtop}{0.5mm}        % gap: box top -> first line of body text
\newcommand{\bookboxpadbottom}{0.5mm}     % gap: last line of body text -> box bottom
\newcommand{\bookboxpadtoptitle}{0.5mm}   % gap: title bar top -> title text
\newcommand{\bookboxpadbottomtitle}{0mm}  % gap: title text -> bottom of title bar


% -----------------------------------------------------------------------
% 12. FLOAT PLACEMENT
% -----------------------------------------------------------------------
% Standard LaTeX float-placement tuning parameters (topnumber, bottomnumber,
% etc.), collected here instead of scattered through config.tex. These
% control how aggressively LaTeX will place figures/tables on a page vs.
% pushing them to a float page. Rarely need changing -- touch these only
% if floats are misbehaving (piling up on float-only pages, or refusing
% to place where expected).
\newcommand{\booktopnumber}{3}            % max floats at TOP of a page
\newcommand{\bookbottomnumber}{2}         % max floats at BOTTOM of a page
\newcommand{\booktotalnumber}{4}          % max floats total on a page
\newcommand{\bookdbltopnumber}{3}         % max two-column floats at top of a page
\newcommand{\booktopfraction}{0.9}        % max page fraction floats may occupy at top
\newcommand{\bookbottomfraction}{0.7}     % max page fraction floats may occupy at bottom
\newcommand{\bookdbltopfraction}{0.9}     % same as above, for two-column floats
\newcommand{\booktextfraction}{0.08}      % min page fraction that MUST remain body text
\newcommand{\bookfloatpagefraction}{0.7}  % min fill fraction for a dedicated float page
\newcommand{\bookdblfloatpagefraction}{0.7} % same as above, for two-column float pages


% =========================================================================
% 13. BOOK COVER
% =========================================================================
% The outer book cover (not to be confused with a chapter cover, Section 6
% above). Full-bleed photo, editor logo, title block, and author line.

% --- Cover photo crop anchor ---------------------------------------------
% Controls which part of the source image is kept when it's cropped to
% fit the cover. 0 = show the low edge, 1 = show the high edge, 0.5 =
% centred. Override per-cover via \makebookcover's optional arguments if
% you need a one-off crop different from this default.
\newcommand{\bookcoverphotoxanchor}{0.6}
\newcommand{\bookcoverphotoyanchor}{0.5}

% --- Editor logo box ------------------------------------------------
\newcommand{\bookcoverlogopath}{./figures/cover/logo}  % path WITHOUT extension --
                                                         % the image loader tries
                                                         % pdf/png/jpeg/jpg in turn
\newcommand{\bookcoverlogowidth}{50mm}
\newcommand{\bookcoverlogoheight}{14mm}
\newcommand{\bookcoverlogox}{22mm}        % distance from the LEFT page edge
\newcommand{\bookcoverlogotopgap}{250mm}  % distance from the TOP page edge

% --- Title / subtitle (left-aligned block) -----------------------------
\newcommand{\bookcovertitlefont}{\fontsize{60}{70}\selectfont}
\newcommand{\bookcoversubtitlefont}{\fontsize{25}{30}\selectfont}
\newcommand{\bookcovertitlecolor}{white}
\newcommand{\bookcoversubtitlecolor}{white}
\newcommand{\bookcovertextleft}{22mm}         % left edge of the text block, from page left
\newcommand{\bookcovertextwidth}{150mm}       % wrap width for title/subtitle
\newcommand{\bookcovertitlebottomgap}{180mm}  % distance from the BOTTOM page edge
                                               % (anchored from bottom so the block stays
                                               % put even if title length changes)
\newcommand{\bookcovertitlesubtitlegap}{6mm}  % gap: title -> subtitle

% --- Authors (bottom-right) ---------------------------------------------
\newcommand{\bookcoverauthorfont}{\fontsize{20}{25}\selectfont}
\newcommand{\bookcoverauthorcolor}{white}
\newcommand{\bookcoverauthorright}{22mm}   % distance from the RIGHT page edge
\newcommand{\bookcoverauthorbottom}{50mm}  % distance from the BOTTOM page edge
\newcommand{\bookcoverauthorwidth}{100mm}


% =========================================================================
% 14. COLOPHON / EXPLANATORY PAGE
% =========================================================================
% The "About this document" page (compile timestamp, source link,
% technical team). Font-leak-safe: \makecolophon wraps its font/colour
% changes in \begingroup...\endgroup, so nothing here needs isolating.

\newcommand{\bookcolophontitle}{About this document}
\newcommand{\bookcolophontitlefont}{\Large}
\newcommand{\bookcolophonbodyfont}{\normalsize}
\newcommand{\bookcolophontopgap}{30mm}    % distance from the TOP page edge to the title
\newcommand{\bookcolophontitlegap}{10mm}  % gap: title -> first paragraph
\newcommand{\bookcolophonparagap}{6mm}    % gap BETWEEN paragraphs
\ifbooklinenumbers\AtBeginDocument{\linenumbers}\fi


% =========================================================================
% 1. PAGE GEOMETRY
% =========================================================================
\setlength{\columnsep}{\bookcolsep}

% includehead: "top" measures to the TOP OF THE HEADER, so the header sits
% inside the top margin: blank -> header box -> headsep -> first body line.
\geometry{
  paperwidth  = \bookpagewidth,
  paperheight = \bookpageheight,
  left        = \booklatmargin,
  right       = \booklatmargin,
  top         = \booktopblank,
  headheight  = \bookheadheight,
  headsep     = \bookheadsep,
  includehead,
  bottom      = \bookbottommargin,
  footskip    = \bookfootskip,
  twoside,
}

\setlength{\parindent}{\bookparindent}
\linespread{\booklinespread}

% Equation spacing -- both the "normal" and "short line before/after" variants,
% or LaTeX silently falls back to its own defaults whenever the line before
% or after a display is short.
\AtBeginDocument{%
  \setlength{\abovedisplayskip}{\bookeqspaceabove}%
  \setlength{\abovedisplayshortskip}{\bookeqspaceabove}%
  \setlength{\belowdisplayskip}{\bookeqspacebelow}%
  \setlength{\belowdisplayshortskip}{\bookeqspacebelow}%
}

\makeatletter
\newcommand{\bookkilllastparskip}{\ifdim\lastskip>\z@\vskip-\lastskip\fi}

\let\bookorigequation\equation
\renewcommand{\equation}{\bookkilllastparskip\bookorigequation}

\let\bookorigalign\align
\renewcommand{\align}{\bookkilllastparskip\bookorigalign}

\let\bookoriggather\gather
\renewcommand{\gather}{\bookkilllastparskip\bookoriggather}

\let\bookorigmultline\multline
\renewcommand{\multline}{\bookkilllastparskip\bookorigmultline}

\let\bookorigsubequations\subequations
\renewcommand{\subequations}{\bookkilllastparskip\bookorigsubequations}
\makeatother


% =========================================================================
% 2. BOX-COLOUR LOGIC  (colours themselves are in parameters.tex, sec. 2)
% =========================================================================
\newcommand{\boxcolor}{%
  \ifnum\value{chapter}<\bookpartoneend
    \colorlet{boxcolorcurrent}{boxcolorpart1}%
  \else\ifnum\value{chapter}<\bookparttwoend
    \colorlet{boxcolorcurrent}{boxcolorpart2}%
  \else
    \colorlet{boxcolorcurrent}{boxcolorpart3}%
  \fi\fi
}
\colorlet{boxcolorcurrent}{boxcolorpart1} % safe default before first \boxcolor call


% =========================================================================
% 3. HEADINGS, PARAGRAPHS & FLOAT SPACING
% =========================================================================
\setcounter{secnumdepth}{\booksecnumdepth}
\renewcommand{\thesubsubsection}{(\alph{subsubsection})}
\AtBeginDocument{\setlength{\parskip}{\bookparskip}}

% Chapter heading is not auto-printed (see \startchapter, sec. 5) -- the
% visible heading is hand-built by \chaptercover / \chapterfirstpage.
\titleformat{\chapter}[block]{\bfseries\fontfamily{phv}\selectfont\Large\color{headingcolor}}{}{0pt}{\raggedright\hyphenpenalty=10000\exhyphenpenalty=10000}

\titleformat{\section}[block]
  {\bfseries\fontfamily{phv}\selectfont\large\color{headingcolor}}
  {\makebox[\booksecnumwidth][l]{\thesection}}
  {0pt}{\raggedright\hyphenpenalty=10000\exhyphenpenalty=10000}
\titlespacing*{\section}{0pt}{\bookspacebeforesection}{\bookspaceaftersection}

\titleformat{\subsection}[block]
  {\bfseries\fontfamily{phv}\selectfont\normalsize\color{headingcolor}}
  {\makebox[\booksubsecnumwidth][l]{\thesubsection}}
  {0pt}{\raggedright\hyphenpenalty=10000\exhyphenpenalty=10000}
\titlespacing*{\subsection}{0pt}{\bookspacebeforesubsection}{\bookspaceaftersubsection}

\titleformat{\subsubsection}[block]
  {\bfseries\fontfamily{phv}\selectfont\small\color{headingcolor}}
  {\makebox[\booksubsubsecnumwidth][l]{\thesubsubsection}}
  {0pt}{\raggedright\hyphenpenalty=10000\exhyphenpenalty=10000}
\titlespacing*{\subsubsection}{0pt}{\bookspacebeforesubsubsection}{\bookspaceaftersubsubsection}

\titleformat{\paragraph}[runin]
  {\bfseries\fontfamily{phv}\selectfont\small\color{maintextcolor}}
  {}{0pt}{}
\titlespacing*{\paragraph}{0pt}{\bookspacebeforeparagraph}{\bookspaceafterparagraph}

\AtBeginDocument{%
  \setlength{\floatsep}{\bookfloatsep}%
  \setlength{\textfloatsep}{\booktextfloatsep}%
  \setlength{\intextsep}{\bookintextsep}%
  \setlength{\dblfloatsep}{\bookfloatsep}%
  \setlength{\dbltextfloatsep}{\booktextfloatsep}%
}

\color{maintextcolor}
\hypersetup{
  colorlinks = true,
  linkcolor  = maintextcolor,
  citecolor  = blue!60!black,
  urlcolor   = blue!60!black,
}


% =========================================================================
% 4. PAGE STYLE / HEADERS  (header only, no footer)
% -------------------------------------------------------------------------
% Even page:  page num | Chapter N  short name        (outer = left)
% Odd page:   section num  title | page num           (outer = right)
% =========================================================================
\newcommand{\bookheadpage}[1]{\sffamily\scriptsize\bfseries\color{headingcolor}#1}

% Captures section number+title together (avoids \rightmark's uppercasing).
\makeatletter
\newcommand{\booksectionmarkstore}{}
\renewcommand{\sectionmark}[1]{%
  \renewcommand{\booksectionmarkstore}{\thesection\ \ #1}%
  \markright{#1}%
}
\makeatother

\fancypagestyle{bookheader}{%
  \fancyhf{}%
  \fancyhead[LO]{%
    \sffamily\scriptsize\color{headingcolor}%
    \bookheadpage{\thepage}\quad\textbar\quad\mdseries Chapter~\thechapter\ \ \bookchaptershortstore%
  }%
  \fancyhead[RE]{%
    \sffamily\scriptsize\color{headingcolor}%
    \booksectionmarkstore\quad\textbar\quad\bookheadpage{\thepage}%
  }%
  \renewcommand{\headrulewidth}{\headerrulethickness}%
  \renewcommand{\headrule}{%
    \hspace*{-\booklatmargin}%
    {\color{headerrulecolor}\rule[0pt]{\bookpagewidth}{\headerrulethickness}}%
  }%
  \renewcommand{\footrulewidth}{0pt}%
}
\pagestyle{bookheader}

% Cover pages: no header/rule/page number.
\fancypagestyle{bookcover}{%
  \fancyhf{}%
  \renewcommand{\headrulewidth}{0pt}%
  \renewcommand{\headrule}{}%
  \renewcommand{\footrulewidth}{0pt}%
}

\makeatletter
\newcommand{\bookchaptershortstore}{}
\newcommand{\chaptershortname}[1]{\renewcommand{\bookchaptershortstore}{#1}}
\newcommand{\bookchaptertitlestore}{}
\makeatother


% =========================================================================
% 5. CHAPTER OPENING  (parity toggle in parameters.tex, sec. 5)
% =========================================================================
\makeatletter
% Text shown on an inserted parity-filler page (a page that exists solely
% to push the next chapter onto the correct recto/verso side). Edit the
% wording directly here if you want something different.
\newcommand{\bookblankpagetext}{Intentionally blank page for page scheme}
\newcommand{\shipbookblankpage}{%
  \onecolumn
  \thispagestyle{empty}%
  \vspace*{\fill}
  \begin{center}
    \sffamily\itshape\footnotesize\color{gray!50}\bookblankpagetext
  \end{center}
  \vspace*{\fill}
  \newpage
  \twocolumn
}
\newcommand{\cleartoevenpage}{%
  \clearpage
  \if@twoside
    \ifodd\c@page\shipbookblankpage\fi
  \fi
}
\newcommand{\cleartooddpage}{%
  \clearpage
  \if@twoside
    \ifodd\c@page\else\shipbookblankpage\fi
  \fi
}
\newcommand{\clearchapterpage}{\ifchapterstartseven\cleartoevenpage\else\cleartooddpage\fi}
\newcommand{\clearopeningpage}{\ifchapterstartseven\cleartooddpage\else\cleartoevenpage\fi}
\makeatother

% Registers the chapter (counter/ToC/marks/counters) without printing
% anything -- the visible heading is built by \chaptercover / \chapterfirstpage.
\makeatletter
\newcommand{\startchapter}[1]{%
  \clearchapterpage
  \refstepcounter{chapter}%
  \renewcommand{\bookchaptertitlestore}{#1}%
  \renewcommand{\bookchaptershortstore}{#1}% default; override with \chaptershortname
  \addcontentsline{toc}{chapter}{\protect\numberline{\thechapter}#1}%
  \markboth{#1}{}%
  \boxcolor
  \setcounter{workedexample}{0}%
}
% Preface / unnumbered front matter.
\newcommand{\startchapterstar}[1]{%
  \clearopeningpage
  \renewcommand{\bookchaptertitlestore}{#1}%
  \renewcommand{\bookchaptershortstore}{#1}%
  \phantomsection
  \addcontentsline{toc}{chapter}{#1}%
  \markboth{#1}{}%
  \boxcolor
}
\makeatother




% =========================================================================
% 6. CHAPTER COVER PAGE
% =========================================================================
% Full-bleed photo, tier fallback T1 -> T0, ext fallback pdf/png/jpeg/jpg:
%   ./figures/chapter{NN}/T{1,0}/{label}_T{1,0}.(pdf|png|jpeg|jpg)
% Chapter number extracted via \xdef (global, fully expanded) to avoid the
% scoping bug where a local \StrLeft/\StrRight result vanished outside its
% group; every path is written out in full (no shared base-path macro) so
% nothing collides if a chapter file redefines a same-named macro.
\makeatletter
\newcommand{\makecoverimage}[1]{%
  \StrLeft{#1}{3}[\covtmpA]\StrRight{\covtmpA}{2}[\covtmpB]%
  \xdef\bookcovchapnum{\covtmpB}%
  \AddToShipoutPictureBG*{%
  \AtPageUpperLeft{\raisebox{-\bookcoverimageheight}{%
    \IfFileExists{./figures/chapter\bookcovchapnum/T1/#1_T1.pdf}{%
      \includegraphics[width=\bookpagewidth,height=\bookcoverimageheight]{./figures/chapter\bookcovchapnum/T1/#1_T1.pdf}%
    }{\IfFileExists{./figures/chapter\bookcovchapnum/T1/#1_T1.png}{%
      \includegraphics[width=\bookpagewidth,height=\bookcoverimageheight]{./figures/chapter\bookcovchapnum/T1/#1_T1.png}%
    }{\IfFileExists{./figures/chapter\bookcovchapnum/T1/#1_T1.jpeg}{%
      \includegraphics[width=\bookpagewidth,height=\bookcoverimageheight]{./figures/chapter\bookcovchapnum/T1/#1_T1.jpeg}%
    }{\IfFileExists{./figures/chapter\bookcovchapnum/T1/#1_T1.jpg}{%
      \includegraphics[width=\bookpagewidth,height=\bookcoverimageheight]{./figures/chapter\bookcovchapnum/T1/#1_T1.jpg}%
    }{\IfFileExists{./figures/chapter\bookcovchapnum/T0/#1_T0.pdf}{%
      \includegraphics[width=\bookpagewidth,height=\bookcoverimageheight]{./figures/chapter\bookcovchapnum/T0/#1_T0.pdf}%
    }{\IfFileExists{./figures/chapter\bookcovchapnum/T0/#1_T0.png}{%
      \includegraphics[width=\bookpagewidth,height=\bookcoverimageheight]{./figures/chapter\bookcovchapnum/T0/#1_T0.png}%
    }{\IfFileExists{./figures/chapter\bookcovchapnum/T0/#1_T0.jpeg}{%
      \includegraphics[width=\bookpagewidth,height=\bookcoverimageheight]{./figures/chapter\bookcovchapnum/T0/#1_T0.jpeg}%
    }{\IfFileExists{./figures/chapter\bookcovchapnum/T0/#1_T0.jpg}{%
      \includegraphics[width=\bookpagewidth,height=\bookcoverimageheight]{./figures/chapter\bookcovchapnum/T0/#1_T0.jpg}%
    }{%
      \colorbox{gray!25}{\parbox[b][\bookcoverimageheight][c]{\bookpagewidth}{%
        \centering\sffamily\large\color{gray!60}[cover photo: #1]}}%
    }}}}}}}}%
  }}%
}%
  % Reserve the part of the image that intrudes into the normal text block.
  \rule{0pt}{\dimexpr\bookcoverimageheight-\booktopmargintotal\relax}\par
}
\makeatother

% Number/caption bar under the photo. Full \bookpagewidth background, but
% content constrained to \booktextwidth. Both cells share one background
% colour (\boxcolorcurrent) so the bar reads as one continuous band.
\newsavebox{\bookcovercapbox}
\newlength{\bookcoverbarht}

\newcommand{\makecoverbar}[1]{%
  \sbox{\bookcovercapbox}{%
    \parbox[c][][c]{\bookcovercapwidth}{%
      \noindent\setlength{\parindent}{0pt}%
      \linespread{\bookcovercapspread}\selectfont
      \sffamily\footnotesize\color{bookcovercapcolor}\justifying #1}}%
  \setlength{\bookcoverbarht}{\dimexpr\ht\bookcovercapbox+\dp\bookcovercapbox+\bookcoverbarpad\relax}%
  \ifdim\bookcoverbarht<\bookcoverbarheight
    \setlength{\bookcoverbarht}{\bookcoverbarheight}%
  \fi
  \noindent\hspace*{-\booklatmargin}%
  \colorbox{boxcolorcurrent}{%
    \parbox[c][\bookcoverbarht][c]{\booklatmargin}{\hfill}%
    \parbox[c][\bookcoverbarht][c]{\dimexpr\booktextwidth-\bookcovercapwidth\relax}{%
      \centering\sffamily\bookcovernumfont\color{bookcovernumcolor}\thechapter}%
    \parbox[c][\bookcoverbarht][c]{\bookcovercapwidth}{\usebox{\bookcovercapbox}}%
    \parbox[c][\bookcoverbarht][c]{\booklatmargin}{\hfill}%
  }%
  \par
}

% "In this chapter ..." label -- standalone, width-constrained to column 1.
\newcommand{\bookcoverlabel}{%
  \parbox{\bookcolwidth}{\noindent\textit{In this chapter} \ldots}%
}

\newcommand{\inthischapter}[1]{%
\noindent #1

}

% Label + gap + two-column body, inside a fixed-height container.
\newcommand{\makecoverspoiler}[1]{%
  \parbox[b][\bookcoverspoilerht][b]{\booktextwidth}{%
    \bookcoverlabel

    \vspace{\bookcoverlabelgap}

    \begin{multicols}{2}
      \inthischapter{#1}%
    \end{multicols}
  }%
}
\newlength{\bookcoverspoilerht}

\newcommand{\chaptercover}[3]{%
  \twocolumn[{%
    \thispagestyle{bookcover}%
    \makecoverimage{#1}%
    \makecoverbar{#2}%
    \vspace{\bookcovergaptobar}%
    \setlength{\bookcoverspoilerht}{\dimexpr
      \bookpageheight-\bookbottommargin-\bookcoverimageheight
      -\bookcoverbarht-\bookcovergaptobar\relax}%
    \makecoverspoiler{#3}%
  }]%
  \thispagestyle{bookcover}%
  \clearpage
}


% =========================================================================
% 7. CHAPTER FIRST PAGE  (title + number box top; skills box top of col. 2)
% =========================================================================
% Header block carries only the title row; column 1 starts right beneath
% it. The skills box is an ordinary float pinned to column 2's top:
% \suppressfloats[t] clears column 1's top only, so column 2 is still free.
% #1 = professional-skills content. Ignored for chapters below
% \bookfirstdecoratedchapter (e.g. the Preface, chapter 0) -- pass {} there.
\newcommand{\chapterfirstpage}[1]{%
  \twocolumn[{%
    \thispagestyle{bookheader}%
    \vspace*{-\dimexpr\bookheadsep-\bookfirstpagetitlenudge\relax}%
    \noindent
    \begin{minipage}[c]{\dimexpr\booktextwidth-\bookfirstpagenumboxwidth-\bookfirstpagetitlenumgap\relax}
      \raggedright\bookfirstpagetitlesize\bfseries\fontfamily{phv}\selectfont
      \color{headingcolor}\bookchaptertitlestore\par
    \end{minipage}%
    \hfill
    \colorbox{boxcolorcurrent}{\parbox[c][\bookfirstpagenumboxheight][c]{\bookfirstpagenumboxwidth}{%
      \centering\sffamily\bookcovernumfont\color{bookcovernumcolor}\thechapter}}%
    \par\vspace{\bookfirstpagetitlegap}%
  }]%
  \thispagestyle{bookheader}%
  \ifnum\value{chapter}<\bookfirstdecoratedchapter\else
    \suppressfloats[t]%
    \begin{figure}[t]
      \begin{boxskills}
      #1
      \end{boxskills}
    \end{figure}%
  \fi
}

% =========================================================================
% 8. ENDING PAGE / HOMEWORK PAGES
% =========================================================================
% Take-home box always spans both columns, bottom of page (invariant) --
% baked into the environment so chapter files need no float wrapper.
\newtcolorbox{boxmessagesinner}{
  myboxstyle, width=\linewidth,
  title = {Take home messages},
}
\newenvironment{boxmessages}
  {\onecolumn\thispagestyle{bookheader}\begin{boxmessagesinner}}
  {\end{boxmessagesinner}\vspace*{\fill}\twocolumn}

% \newpage only breaks columns in twocolumn mode; \clearpage guarantees a
% fresh page, needed before the take-home box and before "Homework".
\newcommand{\chapterhardbreak}{\clearpage}


% =========================================================================
% 9. FIGURES  (S = one column, L = two columns, M = 126mm + side caption)
% =========================================================================

% Centre-crop loader: fills a FIXED width x height box completely, cropping
% overflow evenly from the centre (vs. \bookfitimage's letterboxing).
\newsavebox{\bookcropimgbox}
\newlength{\bookcropimgnw}
\newlength{\bookcropimgnh}
\newcommand{\bookcropimage}[3]{%
  % #1 = image path  #2 = target width  #3 = target height
  \sbox{\bookcropimgbox}{\includegraphics{#1}}%
  \setlength{\bookcropimgnw}{\wd\bookcropimgbox}%
  \setlength{\bookcropimgnh}{\ht\bookcropimgbox}%
  \begin{tikzpicture}
    \path[use as bounding box] (0,0) rectangle (#2,#3);
    \clip (0,0) rectangle (#2,#3);
    \pgfmathtruncatemacro{\bookcropimgwider}{%
      (#2/#3) > (\bookcropimgnw/\bookcropimgnh) ? 1 : 0}%
    \ifnum\bookcropimgwider=1
      \node[inner sep=0pt] at (0.5*#2,0.5*#3) {\includegraphics[width=#2]{#1}};
    \else
      \node[inner sep=0pt] at (0.5*#2,0.5*#3) {\includegraphics[height=#3]{#1}};
    \fi
  \end{tikzpicture}%
}

% Tier/extension fallback for the crop loader (same search order as below).
\newcommand{\bookincludeimagecrop}[5]{%
  \bookincludeimagecropbase{#1}{#3}{#4}{#5}{#2}%
}
\newcommand{\bookincludeimagecropbase}[5]{%
  \IfFileExists{#1.pdf}{\bookcropimage{#1.pdf}{#3}{#4}}{%
  \IfFileExists{#1.png}{\bookcropimage{#1.png}{#3}{#4}}{%
  \IfFileExists{#1.jpeg}{\bookcropimage{#1.jpeg}{#3}{#4}}{%
  \IfFileExists{#1.jpg}{\bookcropimage{#1.jpg}{#3}{#4}}{%
    \IfStrEq{#5}{}{%
      \bookimageplaceholder{#2}{#3}{#4}{portrait}%
    }{%
      \bookincludeimagecropbase{#5}{#2}{#3}{#4}{}%
    }%
  }}}}%
}

% Unified loader: TIER fallback (T1 -> T0), then EXTENSION fallback
% (pdf -> png -> jpeg -> jpg). #1/#2 = T1/T0 base paths, #3 = placeholder
% label, #4/#5 = target width/height, #6 = human-readable size label.
\newcommand{\bookincludeimage}[6]{%
  \bookincludeimagebase{#1}{#3}{#4}{#5}{#6}{#2}%
}
\newcommand{\bookincludeimagebase}[6]{%
  \IfFileExists{#1.pdf}{\bookfitimage{#1.pdf}{#3}{#4}}{%
  \IfFileExists{#1.png}{\bookfitimage{#1.png}{#3}{#4}}{%
  \IfFileExists{#1.jpeg}{\bookfitimage{#1.jpeg}{#3}{#4}}{%
  \IfFileExists{#1.jpg}{\bookfitimage{#1.jpg}{#3}{#4}}{%
    \IfStrEq{#6}{}{%
      \bookimageplaceholder{#2}{#3}{#4}{#5}%
    }{%
      \bookincludeimagebase{#6}{#2}{#3}{#4}{#5}{}%
    }%
  }}}}%
}

% Once a real file is found, only WIDTH is fixed; height follows the
% image's own aspect ratio. #3 (target height) is used only by the
% placeholder fallback below.
\newcommand{\bookfitimage}[3]{%
  \parbox[t]{#2}{\centering\includegraphics[width=#2,keepaspectratio]{#1}}%
}
\newcommand{\bookimageplaceholder}[4]{%
  \fbox{\parbox[c][\dimexpr#3-2\fboxsep-2\fboxrule\relax][c]{%
        \dimexpr#2-2\fboxsep-2\fboxrule\relax}{%
    \centering\sffamily\footnotesize\color{gray!70}
    Figure\\#1\\[2mm]{\tiny #4}}}%
}

% Main figure command -- same name as the manuscript, extended with two
% leading optional args: size (S/M/L, default L) and placement (t/b,
% default t). Existing \figplaceholder{label}{name}{caption} calls keep
% working unchanged (they get the L/t default). Chapter number extraction
% uses \xdef for the same reason as \makecoverimage.
\NewDocumentCommand{\figplaceholder}{O{L} O{t} m m m}{%
  \StrLeft{#3}{3}[\figtmpA]\StrRight{\figtmpA}{2}[\figtmpB]%
  \xdef\bookfigchapnum{\figtmpB}%
  \IfStrEq{#1}{S}{%
    \begin{figure}[#2]
    \centering
    \bookincludeimage
      {./figures/chapter\bookfigchapnum/T1/#3_T1}%
      {./figures/chapter\bookfigchapnum/T0/#3_T0}%
      {#4}{\bookfigwidthS}{\bookfigheightS}{\bookfiglabelS}
    \captionsetup{aboveskip=\bookfigcapabove, belowskip=\bookfigcapbelow}%
    \caption{#5}\label{#3}
    \end{figure}%
  }{%
  \IfStrEq{#1}{M}{%
    \begin{figure*}[#2]
    \noindent
    % Both minipages [t] AND both opening with \vspace{0pt}: the \vspace
    % is what fixes a minipage whose first item is a box (a \parbox[c])
    % from taking that box's CENTRE as its baseline reference, which
    % would otherwise pull the caption down to mid-image height.
    \begin{minipage}[t]{\bookMfigwidth}
      \vspace{0pt}%
      \centering\bookincludeimage
        {./figures/chapter\bookfigchapnum/T1/#3_T1}%
        {./figures/chapter\bookfigchapnum/T0/#3_T0}%
        {#4}{\linewidth}{\bookfigheightM}{\bookfiglabelM}
    \end{minipage}\hfill%
    \begin{minipage}[t]{\dimexpr\booktextwidth-\bookMfigwidth-\bookcolsep\relax}
      \vspace{0pt}%
      \captionof{figure}{#5}\label{#3}
    \end{minipage}
    \end{figure*}%
  }{%
    \begin{figure*}[#2]
    \centering
    \bookincludeimage
      {./figures/chapter\bookfigchapnum/T1/#3_T1}%
      {./figures/chapter\bookfigchapnum/T0/#3_T0}%
      {#4}{\bookfigwidthL}{\bookfigheightL}{\bookfiglabelL}
    \captionsetup{aboveskip=\bookfigcapabove, belowskip=\bookfigcapbelow}%
    \caption{#5}\label{#3}
    \end{figure*}%
  }}%
}

% Width-only loader/fallback: image fills a target WIDTH, height follows
% naturally (no letterboxing). Used by \boxfigplaceholder.
\newcommand{\bookfitimagewidth}[2]{%
  \parbox{#2}{\centering\includegraphics[width=#2]{#1}}%
}
\newcommand{\bookincludeimagewidth}[4]{%
  \bookincludeimagewidthbase{#1}{#3}{#4}{#2}%
}
\newcommand{\bookincludeimagewidthbase}[4]{%
  \IfFileExists{#1.pdf}{\bookfitimagewidth{#1.pdf}{#3}}{%
  \IfFileExists{#1.png}{\bookfitimagewidth{#1.png}{#3}}{%
  \IfFileExists{#1.jpeg}{\bookfitimagewidth{#1.jpeg}{#3}}{%
  \IfFileExists{#1.jpg}{\bookfitimagewidth{#1.jpg}{#3}}{%
    \IfStrEq{#4}{}{%
      \bookimageplaceholder{#2}{#3}{\bookboxfigheight}{box image}%
    }{%
      \bookincludeimagewidthbase{#4}{#2}{#3}{}%
    }%
  }}}}%
}

\newcommand{\boxfigplaceholder}[1]{%
  \StrLeft{#1}{3}[\figtmpA]\StrRight{\figtmpA}{2}[\figtmpB]%
  \xdef\bookfigchapnum{\figtmpB}%
  \begin{center}
  \bookincludeimagewidth
    {./figures/chapter\bookfigchapnum/T1/#1_T1}%
    {./figures/chapter\bookfigchapnum/T0/#1_T0}%
    {#1}{\linewidth}%
  \end{center}
}

\newcommand{\boxbiofigplaceholder}[1]{%
  \StrLeft{#1}{3}[\figtmpA]\StrRight{\figtmpA}{2}[\figtmpB]%
  \xdef\bookfigchapnum{\figtmpB}%
  \centering
  \bookincludeimagecrop
    {./figures/chapter\bookfigchapnum/T1/#1_T1}%
    {./figures/chapter\bookfigchapnum/T0/#1_T0}%
    {#1}{\bookbiofigwidth}{\bookbiofigheight}%
}


% =========================================================================
% 10. BOXES  (S/L sized via the call site -- see transform_manuscript.py)
% =========================================================================
\tcbset{
  myboxstyle/.style={
    enhanced,
    % NOT breakable: every box lives inside a float, which cannot break
    % across pages -- a breakable box inside one renders corrupted. A box
    % that won't fit on one page must be split editorially.
    breakable    = false,
    colback      = boxcolorcurrent,
    colbacktitle = boxcolorcurrent,
    coltitle     = headingcolor,
    colframe     = boxcolorcurrent,
    boxrule      = 0pt,
    arc          = 0pt,
    boxsep       = \bookboxpad,
    left         = \bookboxpadlat,
    right        = \bookboxpadlat,
    top          = \bookboxpadtop,
    bottom       = \bookboxpadbottom,
    toptitle     = \bookboxpadtoptitle,
    bottomtitle  = \bookboxpadbottomtitle,
    fonttitle    = \bfseries\sffamily\small,
    fontupper    = \sffamily\footnotesize\color{maintextcolor},
    fontlower    = \sffamily\footnotesize\color{maintextcolor},
  }
}

% Numbered box -- same keys/theme/title signature as the manuscript's
% boxmain, renamed boxmaininner. \boxmain (below) owns size/placement.
\newtcolorbox[
  auto counter,
  number within = chapter,
  list inside   = lob,
]{boxmaininner}[3][]{
  myboxstyle, width=\linewidth,
  title = {Box~\thetcbcounter\quad#2: #3},
  #1,
}
\crefname{tcb@cnt@boxmaininner}{box}{boxes}
\Crefname{tcb@cnt@boxmaininner}{Box}{Boxes}

% \ifbookboxsizeS records which literal float (figure/figure*) THIS
% \boxmain call opened, so the matching \end{...} closes the same one.
\NewDocumentEnvironment{boxmain}{O{L} O{t} O{} m m}{%
  \IfStrEq{#1}{S}{\bookboxsizeStrue}{\bookboxsizeSfalse}%
  \ifbookboxsizeS
    \begin{figure}[#2]
  \else
    \begin{figure*}[#2]
  \fi
  \begin{boxmaininner}[#3]{#4}{#5}%
}{%
  \end{boxmaininner}
  \ifbookboxsizeS
    \end{figure}
  \else
    \end{figure*}
  \fi
}

% Professional-skills box -- placed by \chapterfirstpage.
\newtcolorbox{boxskills}{
  myboxstyle,
  title = {Professional practice relevance and skills developed},
}

% Biographies -- always two-column width, bottom of page (invariant), so
% the float wrapper is baked in; chapter files call it as in the manuscript.
\newtcolorbox{boxbioinner}[3]{
  myboxstyle, width=\linewidth,
  breakable=false, sidebyside, sidebyside align=top,
  lefthand width=\bookbiofigwidth,
  sidebyside gap=\bookbiogap,
  before lower={%
    {\bfseries\sffamily\small\color{headingcolor}Biography: #1 (#2 -- #3)}\par
    \vspace{\bookbiotitlegap}%
  },
}
% Optional leading [t/b] controls page placement, default b per spec.
\newenvironment{boxbio}[4][b]
  {\begin{figure*}[#1]\begin{boxbioinner}{#2}{#3}{#4}}
  {\end{boxbioinner}\end{figure*}}


% =========================================================================
% 11. WORKED EXAMPLES  (inline, not floated)
% =========================================================================
\newcounter{workedexample}
\numberwithin{workedexample}{chapter}
\renewcommand{\theworkedexample}{\thechapter.\arabic{workedexample}}

\newenvironment{workedexample}[1]
{%
  \refstepcounter{workedexample}%
  \setlength{\parindent}{0pt}%
  \par\medskip\noindent
  {\bfseries\small\fontfamily{phv}\selectfont\color{headingcolor}%
   Worked Example~\theworkedexample: #1}\par
}
{%
  \par\hfill End of Worked Example $\blacksquare$\par\vspace{\bookspacebeforeexample}
}
\newcommand{\subcase}[1]{\par\noindent\textbf{#1.}\enspace\ignorespaces}
\crefname{workedexample}{Worked Example}{Worked Examples}
\Crefname{workedexample}{Worked Example}{Worked Examples}

\numberwithin{equation}{chapter}
\numberwithin{figure}{chapter}
\numberwithin{table}{chapter}


% =========================================================================
% 12. FLOAT PLACEMENT
% =========================================================================
% Figure-dense book: floats may only sit at top/bottom of a page, so LaTeX
% needs to be allowed to fill much more of a page than the defaults permit.
\setcounter{topnumber}{\booktopnumber}
\setcounter{bottomnumber}{\bookbottomnumber}
\setcounter{totalnumber}{\booktotalnumber}
\setcounter{dbltopnumber}{\bookdbltopnumber}
\renewcommand{\topfraction}{\booktopfraction}
\renewcommand{\bottomfraction}{\bookbottomfraction}
\renewcommand{\dbltopfraction}{\bookdbltopfraction}
\renewcommand{\textfraction}{\booktextfraction}
\renewcommand{\floatpagefraction}{\bookfloatpagefraction}
\renewcommand{\dblfloatpagefraction}{\bookdblfloatpagefraction}